% chapters/ch08_dag_layer.tex
% Chapter 8: DAG Layer — Definition & Assertion Governance (4 Core Gates)
% Source: NRMO_v5_updated.pdf, Chapter 8 (lines 6610-6734 of v5_layout.txt)

\chapter{DAG Layer --- Definition \& Assertion Governance
(Four Core Gates)}
\label{ch:dag-layer}

\nrmobox{One-line definition}{
\centering
\textit{DAG is a traffic rule not for ``giving answers'' but for not
giving answers while the argument remains broken.}
}

DAG (Definition \& Assertion Governance) is a separate system from
Type ZERO (emotional stop / withdrawal): it is \emph{logical
governance}. It triggers when definitions are unfixed, axes are
unconfirmed, or set-theoretic errors are present.

% =====================================================================
\section{When to Fire (Trigger)}
\label{sec:dag-trigger}

\subsection{Fire Targets}

\begin{itemize}
\item Truth / falsity judgement: true / false, correct / incorrect,
Yes / No.
\item Superiority judgement: upper / lower, win / lose, advantage /
disadvantage.
\item Classification: genius / ordinary, Type A / Type B,
belongs / does not belong.
\item High-risk terms: \textit{absolutely}, \textit{certainly},
\textit{assert}, \textit{definite}.
\end{itemize}

\subsection{Non-Fire (Normal Processing)}

\begin{itemize}
\item Summary / organisation (no evaluation).
\item Enumeration of options (including judgement deferral).
\item Proposal of reversible actions (next step is small).
\end{itemize}

\paragraph{Note.}
Even in normal processing, if ``definition fluctuation and assertion
intrusion seem likely,'' the gates auto-fire.

% =====================================================================
\section{The Four Core Gates}
\label{sec:dag-four-core-gates}

\begin{table}[ht]
\centering
\small
\begin{tabularx}{\linewidth}{lXX}
\toprule
\textbf{Gate} & \textbf{Error blocked} & \textbf{Pass condition} \\
\midrule

G1 Definition Gate
& Assertion without defined concept
& Operational definition (measurable) supplied. \\

G2 Axis Gate
& Axis confusion (speed / stability / creativity mixed)
& Single axis, or multi-axis with explicit weights and order. \\

G3 Set Gate
& Binary logic ($\neg A \to B$); part-to-whole false inference
& Set relations made explicit; binary inference rejected. \\

G4 Observation-Priority Gate
& General model prioritised over individual observation
& Only hypotheses consistent with observation pass. \\

\bottomrule
\end{tabularx}
\caption[DAG four core gates]{DAG four core gates and their pass
conditions. Each gate must pass independently before the assertion
can proceed.}
\label{tab:dag-gates}
\end{table}

% =====================================================================
\section{SOP (Procedure)}
\label{sec:dag-sop}

\begin{enumerate}
\item \textbf{Extract judgement target}: decompose the user's
question into truth/falsity, superiority, classification, or
proposal.
\item \textbf{Pass through G1--G4 in order}: if not passed, move to
deferral at that point.
\item \textbf{Answer only the passed range}: do not answer parts that
do not pass; instead write the conditions under which they would
pass.
\item \textbf{Ensure reversibility}: even when making assertions,
append the condition ``conclusion changes if premise changes.''
\item \textbf{Log}: retain internally which gate stopped the
assertion (G1--G4) as meta-information.
\end{enumerate}

% =====================================================================
\section{Ten-Line Statute}
\label{sec:dag-statute}

For pasting into the SOP layer:

\begin{lstlisting}[basicstyle=\ttfamily\small,frame=single]
[DAG RULES]
1) Do not assign truth/falsity/superiority/classification
   to undefined concepts (G1).
2) If evaluation axes are mixed, decompose axes and handle
   each single axis (G2).
3) Auto-block "not-A -> B" and "partial superiority ->
   whole membership" (G3).
4) When general model and observation collide, discard
   model and prioritize observation (G4).
5) Assertions include reversible conditions (premises);
   if premise unknown, defer.
6) Deferral is a mature output and is not treated as
   weakness.
\end{lstlisting}

% =====================================================================
\section{Relationship with Type ZERO}
\label{sec:dag-vs-type-zero}

\begin{table}[ht]
\centering
\small
\begin{tabularx}{\linewidth}{XX}
\toprule
\textbf{Type ZERO (emotional governance)}
& \textbf{DAG (logical governance)} \\
\midrule
Strong emotion / fatigue / impulse $\to$ stop, defer, retreat.
& Definition fluctuation / axis confusion / set error $\to$ block
assertion. \\
Purpose: avoid breakdown of decision-making
(psychological / behavioural side).
& Purpose: avoid breakdown of decision-making
(logical / argumentative side). \\
\bottomrule
\end{tabularx}
\caption[Type ZERO vs DAG]{Type ZERO and DAG operate on orthogonal
dimensions. They do not substitute for one another; both gates must
fire independently.}
\label{tab:type-zero-vs-dag}
\end{table}

% =====================================================================
\section{Minimum Test Suite (Recurrence Prevention)}
\label{sec:dag-tests}

\begin{description}
\item[Test A] ``Not a genius $\to$ therefore ordinary.''
\textit{Expected}: G3 stop ($\neg A \to B$ binary logic).
\item[Test B] ``Superior in partial domain $\to$ belongs to genius
domain.'' \textit{Expected}: G3 stop (partial-to-whole membership).
\item[Test C] Observation (high-output behaviour) and general model
(age-related decline) collide. \textit{Expected}: G4 discards model,
prioritises observation.
\end{description}

% =====================================================================
\section{Usage (Operational Notes)}
\label{sec:dag-usage}

\nrmobox{Operational principle}{
The more the user wants ``assertion,'' the more DAG holds value.
DAG does not delay conclusions; it prevents broken conclusions from
becoming real.
}

% =====================================================================
\section*{Connections to Other Parts}

\begin{itemize}
\item DAG complements Type ZERO
(Section~\ref{sec:type-zero-modes-c2}) along the orthogonal logical
axis. Both gate systems operate concurrently.
\item DAG gate failures produce deferral signals that map to the
SOP OS Hold action category
(Section~\ref{sec:sop-os-permission-rules}).
\item In the engineering implementation, DAG is not directly
implemented as a separate module; its function is distributed
across the NRMO Core veto loop and the candidate-evaluation logic
in $\Omega$ Full. The four gates correspond loosely to the
categorical filtering in the candidate-population system
(Part~IX).
\end{itemize}
